\chapter{Methods}

Progress was coordinated through biweekly meetings with the product owner at Ascento and the academic advisor, serving 
as informal checkpoints to review results, reprioritize implementation effort, 
and surface blockers early. No formal process model was imposed: given the exploratory nature of the work and the 
tight feedback loop afforded by working within the company, 
a lightweight review cadence was more appropriate than a structured sprint framework. 
Requirements were not formalized in a separate document. Working directly within Ascento's 
engineering context procided sufficient domain knowledge to derive the operational constraints implicitly: 
false negatives are the critical failure mode, as a missed detection on a live patrol has no downstream safety net, 
whereas false positives are surfaced to human reviewers and therefore recoverable. These constraints translate directly 
into the choice of \gls{map@50} as the primary benchmark metric and the explicit minimization of false negative rate as 
the recondary optimization objective, as defined in !!!!!!!!!
\section{Data and Tools}
% Placeholder — describe datasets, software, and instruments used.
All code and configuration are version-controlled via \gls{git}. 
Dataset variants are tracked independently using \gls{dvc}, 
which maintains a reproducible record of each augmentation stage 
without storing large binary files in the repository. 
Dependency management across the pipeline's heterogeneous components is 
handled via \gls{uv}, ensuring reproducible environments across machines. 
\gls{yolo} training runs are tracked in \gls{wandb} with generation 
parameters and synthetic yiels rates logged directly in the codebase 
alongside training metrics, providing a complete record linking 
each dataset variant to its downstram detector performance.  
Hyperparameter search is automated via \gls{optuna} which runs structured 
sweeps over the training configuration and selects the combination 
that maximizes validation \gls{map@50} within a fixed compute budget. 
Training infrastructure was split across two environments determined 
by hardware constraints. All components except \gls{pbe} were trained 
and run locally on a Apple \gls{m1} chip.
\gls{pbe} requires \gls{cuda} and was therefore executed on the 
\gls{hslu} GPUhub using a single Nvidia A16 core. This split introduced no methodological 
inconsistencies but is noted as a reproducibility consideration: results on the \gls{m1} and 
on \gls{cuda} hardware may differ numerically due to floating point implementation 
differences across backends. 
